%%%%%%%%%%%%%%%%%%%%%%%%%%%%%%%%%%%%%%%%%%%%%%%%%%
% Gemeinsamer Stil für Quellcode und Konsolenausgaben
% Einbinden mit %%%%%%%%%%%%%%%%%%%%%%%%%%%%%%%%%%%%%%%%%%%%%%%%%%
% Gemeinsamer Stil für Quellcode und Konsolenausgaben
% Einbinden mit %%%%%%%%%%%%%%%%%%%%%%%%%%%%%%%%%%%%%%%%%%%%%%%%%%
% Gemeinsamer Stil für Quellcode und Konsolenausgaben
% Einbinden mit %%%%%%%%%%%%%%%%%%%%%%%%%%%%%%%%%%%%%%%%%%%%%%%%%%
% Gemeinsamer Stil für Quellcode und Konsolenausgaben
% Einbinden mit \input{code-stil} in der Präambel.
%
% Quellcode (mit Dateiname als Reiter):
%   \begin{javacode}{Pizza.java} ... \end{javacode}
%   \begin{javacode}[Ausschnitt]{Beispiel} ... \end{javacode}   (ohne Zeilennummern)
%   \begin{javacode}[breakable]{Lang.java} ... \end{javacode}  (darf über Seiten umbrechen)
%
% Konsolenausgabe (Eingaben des Benutzers zwischen @...@):
%   \begin{konsole} Bitte Jahr eingeben:
%   @2026@ \end{konsole}
%
% Leere Konsole zum Ausfüllen (z. B. in Tests), Höhe in Zeilen:
%   \leerekonsole[5]
%%%%%%%%%%%%%%%%%%%%%%%%%%%%%%%%%%%%%%%%%%%%%%%%%%

\RequirePackage{xcolor}
\RequirePackage{listings}
\RequirePackage[most]{tcolorbox}
\RequirePackage[varqu,varl,scaled=0.95]{zi4}   % Inconsolata als Schreibmaschinenschrift

% Farben passend zu den Grafiken (Syntaxdiagramme, Dreieck, Pizza)
\definecolor{diagramm}{RGB}{74,144,217}
\definecolor{codeHintergrund}{RGB}{244,248,253}
\definecolor{codeSchluessel}{RGB}{28,86,160}
\definecolor{codeKommentar}{RGB}{125,132,145}
\definecolor{codeString}{RGB}{40,128,72}
\definecolor{codeZeilennr}{RGB}{150,170,195}
\definecolor{konsoleHintergrund}{RGB}{238,239,242}
\definecolor{konsoleText}{RGB}{45,50,60}
\definecolor{konsoleEingabe}{RGB}{204,102,0}
\definecolor{konsoleReiter}{RGB}{110,118,130}

% Umlaute in Kommentaren und Strings
\lstdefinestyle{gdiUmlaute}{
	literate=
		{ä}{{\"a}}1 {ö}{{\"o}}1 {ü}{{\"u}}1
		{Ä}{{\"A}}1 {Ö}{{\"O}}1 {Ü}{{\"U}}1
		{ß}{{\ss}}1
}

\lstdefinestyle{gdiJava}{
	style=gdiUmlaute,
	aboveskip=0pt,
	belowskip=0pt,
	language=Java,
	basicstyle=\ttfamily\footnotesize,
	keywordstyle=\color{codeSchluessel}\bfseries,
	commentstyle=\color{codeKommentar}\itshape,
	stringstyle=\color{codeString},
	numbers=left,
	numberstyle=\scriptsize\ttfamily\color{codeZeilennr},
	numbersep=10pt,
	tabsize=4,
	columns=fullflexible,
	keepspaces=true,
	showstringspaces=false,
	breaklines=true,
	postbreak=\space
}

\lstdefinestyle{gdiKonsole}{
	style=gdiUmlaute,
	aboveskip=0pt,
	belowskip=0pt,
	language={},
	basicstyle=\ttfamily\footnotesize\color{konsoleText},
	columns=fullflexible,
	keepspaces=true,
	breaklines=true,
	moredelim=[is][\color{konsoleEingabe}\bfseries]{@}{@}
}

% Gemeinsames Aussehen der Kästen
\tcbset{
	gdiKasten/.style={
		enhanced, breakable,
		listing only,
		sharp corners,
		boxrule=0pt,
		top=1.5mm, bottom=1.5mm, right=2mm,
		before skip=3mm, after skip=3mm,
		attach boxed title to top left={xshift=0mm, yshift=-0.2mm},
		boxed title style={sharp corners, boxrule=0pt, top=0.6mm, bottom=0.6mm, left=2mm, right=2mm},
		fonttitle=\ttfamily\footnotesize\bfseries
	}
}

% Quellcode: heller Kasten mit blauer Linie links und Dateiname als Reiter
\newtcblisting{javacode}[2][]{
	gdiKasten,
	listing options={style=gdiJava},
	colback=codeHintergrund,
	unbreakable,
	borderline west={2pt}{0pt}{diagramm},
	left=8mm,
	title={#2},
	coltitle=white,
	boxed title style={colback=diagramm},
	#1
}
% Kurzer Ausschnitt ohne Zeilennummern
\tcbset{Ausschnitt/.style={listing options={style=gdiJava, numbers=none}, left=3mm}}

% Konsolenausgabe: hellgrauer Kasten (druckerfreundlich), Eingaben orange
\newtcblisting{konsole}[1][Konsole]{
	gdiKasten,
	listing options={style=gdiKonsole},
	colback=konsoleHintergrund,
	unbreakable,
	left=3mm,
	title={#1},
	borderline west={2pt}{0pt}{konsoleReiter},
	coltitle=white,
	boxed title style={colback=konsoleReiter}
}

% Leere Konsole zum Ausfüllen (z. B. in Tests): \leerekonsole[Anzahl Zeilen]
\newcommand{\leerekonsole}[1][5]{%
	\begin{tcolorbox}[
		enhanced, unbreakable, sharp corners, boxrule=0pt,
		colback=konsoleHintergrund,
		borderline west={2pt}{0pt}{konsoleReiter},
		before skip=3mm, after skip=3mm,
		height=\dimexpr #1\dimexpr 8mm\relax + 3mm\relax,
		title={Konsole}, coltitle=white,
		attach boxed title to top left={xshift=0mm, yshift=-0.2mm},
		boxed title style={sharp corners, boxrule=0pt, colback=konsoleReiter, top=0.6mm, bottom=0.6mm, left=2mm, right=2mm},
		fonttitle=\ttfamily\footnotesize\bfseries
	]\end{tcolorbox}%
}
 in der Präambel.
%
% Quellcode (mit Dateiname als Reiter):
%   \begin{javacode}{Pizza.java} ... \end{javacode}
%   \begin{javacode}[Ausschnitt]{Beispiel} ... \end{javacode}   (ohne Zeilennummern)
%   \begin{javacode}[breakable]{Lang.java} ... \end{javacode}  (darf über Seiten umbrechen)
%
% Konsolenausgabe (Eingaben des Benutzers zwischen @...@):
%   \begin{konsole} Bitte Jahr eingeben:
%   @2026@ \end{konsole}
%
% Leere Konsole zum Ausfüllen (z. B. in Tests), Höhe in Zeilen:
%   \leerekonsole[5]
%%%%%%%%%%%%%%%%%%%%%%%%%%%%%%%%%%%%%%%%%%%%%%%%%%

\RequirePackage{xcolor}
\RequirePackage{listings}
\RequirePackage[most]{tcolorbox}
\RequirePackage[varqu,varl,scaled=0.95]{zi4}   % Inconsolata als Schreibmaschinenschrift

% Farben passend zu den Grafiken (Syntaxdiagramme, Dreieck, Pizza)
\definecolor{diagramm}{RGB}{74,144,217}
\definecolor{codeHintergrund}{RGB}{244,248,253}
\definecolor{codeSchluessel}{RGB}{28,86,160}
\definecolor{codeKommentar}{RGB}{125,132,145}
\definecolor{codeString}{RGB}{40,128,72}
\definecolor{codeZeilennr}{RGB}{150,170,195}
\definecolor{konsoleHintergrund}{RGB}{238,239,242}
\definecolor{konsoleText}{RGB}{45,50,60}
\definecolor{konsoleEingabe}{RGB}{204,102,0}
\definecolor{konsoleReiter}{RGB}{110,118,130}

% Umlaute in Kommentaren und Strings
\lstdefinestyle{gdiUmlaute}{
	literate=
		{ä}{{\"a}}1 {ö}{{\"o}}1 {ü}{{\"u}}1
		{Ä}{{\"A}}1 {Ö}{{\"O}}1 {Ü}{{\"U}}1
		{ß}{{\ss}}1
}

\lstdefinestyle{gdiJava}{
	style=gdiUmlaute,
	aboveskip=0pt,
	belowskip=0pt,
	language=Java,
	basicstyle=\ttfamily\footnotesize,
	keywordstyle=\color{codeSchluessel}\bfseries,
	commentstyle=\color{codeKommentar}\itshape,
	stringstyle=\color{codeString},
	numbers=left,
	numberstyle=\scriptsize\ttfamily\color{codeZeilennr},
	numbersep=10pt,
	tabsize=4,
	columns=fullflexible,
	keepspaces=true,
	showstringspaces=false,
	breaklines=true,
	postbreak=\space
}

\lstdefinestyle{gdiKonsole}{
	style=gdiUmlaute,
	aboveskip=0pt,
	belowskip=0pt,
	language={},
	basicstyle=\ttfamily\footnotesize\color{konsoleText},
	columns=fullflexible,
	keepspaces=true,
	breaklines=true,
	moredelim=[is][\color{konsoleEingabe}\bfseries]{@}{@}
}

% Gemeinsames Aussehen der Kästen
\tcbset{
	gdiKasten/.style={
		enhanced, breakable,
		listing only,
		sharp corners,
		boxrule=0pt,
		top=1.5mm, bottom=1.5mm, right=2mm,
		before skip=3mm, after skip=3mm,
		attach boxed title to top left={xshift=0mm, yshift=-0.2mm},
		boxed title style={sharp corners, boxrule=0pt, top=0.6mm, bottom=0.6mm, left=2mm, right=2mm},
		fonttitle=\ttfamily\footnotesize\bfseries
	}
}

% Quellcode: heller Kasten mit blauer Linie links und Dateiname als Reiter
\newtcblisting{javacode}[2][]{
	gdiKasten,
	listing options={style=gdiJava},
	colback=codeHintergrund,
	unbreakable,
	borderline west={2pt}{0pt}{diagramm},
	left=8mm,
	title={#2},
	coltitle=white,
	boxed title style={colback=diagramm},
	#1
}
% Kurzer Ausschnitt ohne Zeilennummern
\tcbset{Ausschnitt/.style={listing options={style=gdiJava, numbers=none}, left=3mm}}

% Konsolenausgabe: hellgrauer Kasten (druckerfreundlich), Eingaben orange
\newtcblisting{konsole}[1][Konsole]{
	gdiKasten,
	listing options={style=gdiKonsole},
	colback=konsoleHintergrund,
	unbreakable,
	left=3mm,
	title={#1},
	borderline west={2pt}{0pt}{konsoleReiter},
	coltitle=white,
	boxed title style={colback=konsoleReiter}
}

% Leere Konsole zum Ausfüllen (z. B. in Tests): \leerekonsole[Anzahl Zeilen]
\newcommand{\leerekonsole}[1][5]{%
	\begin{tcolorbox}[
		enhanced, unbreakable, sharp corners, boxrule=0pt,
		colback=konsoleHintergrund,
		borderline west={2pt}{0pt}{konsoleReiter},
		before skip=3mm, after skip=3mm,
		height=\dimexpr #1\dimexpr 8mm\relax + 3mm\relax,
		title={Konsole}, coltitle=white,
		attach boxed title to top left={xshift=0mm, yshift=-0.2mm},
		boxed title style={sharp corners, boxrule=0pt, colback=konsoleReiter, top=0.6mm, bottom=0.6mm, left=2mm, right=2mm},
		fonttitle=\ttfamily\footnotesize\bfseries
	]\end{tcolorbox}%
}
 in der Präambel.
%
% Quellcode (mit Dateiname als Reiter):
%   \begin{javacode}{Pizza.java} ... \end{javacode}
%   \begin{javacode}[Ausschnitt]{Beispiel} ... \end{javacode}   (ohne Zeilennummern)
%   \begin{javacode}[breakable]{Lang.java} ... \end{javacode}  (darf über Seiten umbrechen)
%
% Konsolenausgabe (Eingaben des Benutzers zwischen @...@):
%   \begin{konsole} Bitte Jahr eingeben:
%   @2026@ \end{konsole}
%
% Leere Konsole zum Ausfüllen (z. B. in Tests), Höhe in Zeilen:
%   \leerekonsole[5]
%%%%%%%%%%%%%%%%%%%%%%%%%%%%%%%%%%%%%%%%%%%%%%%%%%

\RequirePackage{xcolor}
\RequirePackage{listings}
\RequirePackage[most]{tcolorbox}
\RequirePackage[varqu,varl,scaled=0.95]{zi4}   % Inconsolata als Schreibmaschinenschrift

% Farben passend zu den Grafiken (Syntaxdiagramme, Dreieck, Pizza)
\definecolor{diagramm}{RGB}{74,144,217}
\definecolor{codeHintergrund}{RGB}{244,248,253}
\definecolor{codeSchluessel}{RGB}{28,86,160}
\definecolor{codeKommentar}{RGB}{125,132,145}
\definecolor{codeString}{RGB}{40,128,72}
\definecolor{codeZeilennr}{RGB}{150,170,195}
\definecolor{konsoleHintergrund}{RGB}{238,239,242}
\definecolor{konsoleText}{RGB}{45,50,60}
\definecolor{konsoleEingabe}{RGB}{204,102,0}
\definecolor{konsoleReiter}{RGB}{110,118,130}

% Umlaute in Kommentaren und Strings
\lstdefinestyle{gdiUmlaute}{
	literate=
		{ä}{{\"a}}1 {ö}{{\"o}}1 {ü}{{\"u}}1
		{Ä}{{\"A}}1 {Ö}{{\"O}}1 {Ü}{{\"U}}1
		{ß}{{\ss}}1
}

\lstdefinestyle{gdiJava}{
	style=gdiUmlaute,
	aboveskip=0pt,
	belowskip=0pt,
	language=Java,
	basicstyle=\ttfamily\footnotesize,
	keywordstyle=\color{codeSchluessel}\bfseries,
	commentstyle=\color{codeKommentar}\itshape,
	stringstyle=\color{codeString},
	numbers=left,
	numberstyle=\scriptsize\ttfamily\color{codeZeilennr},
	numbersep=10pt,
	tabsize=4,
	columns=fullflexible,
	keepspaces=true,
	showstringspaces=false,
	breaklines=true,
	postbreak=\space
}

\lstdefinestyle{gdiKonsole}{
	style=gdiUmlaute,
	aboveskip=0pt,
	belowskip=0pt,
	language={},
	basicstyle=\ttfamily\footnotesize\color{konsoleText},
	columns=fullflexible,
	keepspaces=true,
	breaklines=true,
	moredelim=[is][\color{konsoleEingabe}\bfseries]{@}{@}
}

% Gemeinsames Aussehen der Kästen
\tcbset{
	gdiKasten/.style={
		enhanced, breakable,
		listing only,
		sharp corners,
		boxrule=0pt,
		top=1.5mm, bottom=1.5mm, right=2mm,
		before skip=3mm, after skip=3mm,
		attach boxed title to top left={xshift=0mm, yshift=-0.2mm},
		boxed title style={sharp corners, boxrule=0pt, top=0.6mm, bottom=0.6mm, left=2mm, right=2mm},
		fonttitle=\ttfamily\footnotesize\bfseries
	}
}

% Quellcode: heller Kasten mit blauer Linie links und Dateiname als Reiter
\newtcblisting{javacode}[2][]{
	gdiKasten,
	listing options={style=gdiJava},
	colback=codeHintergrund,
	unbreakable,
	borderline west={2pt}{0pt}{diagramm},
	left=8mm,
	title={#2},
	coltitle=white,
	boxed title style={colback=diagramm},
	#1
}
% Kurzer Ausschnitt ohne Zeilennummern
\tcbset{Ausschnitt/.style={listing options={style=gdiJava, numbers=none}, left=3mm}}

% Konsolenausgabe: hellgrauer Kasten (druckerfreundlich), Eingaben orange
\newtcblisting{konsole}[1][Konsole]{
	gdiKasten,
	listing options={style=gdiKonsole},
	colback=konsoleHintergrund,
	unbreakable,
	left=3mm,
	title={#1},
	borderline west={2pt}{0pt}{konsoleReiter},
	coltitle=white,
	boxed title style={colback=konsoleReiter}
}

% Leere Konsole zum Ausfüllen (z. B. in Tests): \leerekonsole[Anzahl Zeilen]
\newcommand{\leerekonsole}[1][5]{%
	\begin{tcolorbox}[
		enhanced, unbreakable, sharp corners, boxrule=0pt,
		colback=konsoleHintergrund,
		borderline west={2pt}{0pt}{konsoleReiter},
		before skip=3mm, after skip=3mm,
		height=\dimexpr #1\dimexpr 8mm\relax + 3mm\relax,
		title={Konsole}, coltitle=white,
		attach boxed title to top left={xshift=0mm, yshift=-0.2mm},
		boxed title style={sharp corners, boxrule=0pt, colback=konsoleReiter, top=0.6mm, bottom=0.6mm, left=2mm, right=2mm},
		fonttitle=\ttfamily\footnotesize\bfseries
	]\end{tcolorbox}%
}
 in der Präambel.
%
% Quellcode (mit Dateiname als Reiter):
%   \begin{javacode}{Pizza.java} ... \end{javacode}
%   \begin{javacode}[Ausschnitt]{Beispiel} ... \end{javacode}   (ohne Zeilennummern)
%   \begin{javacode}[breakable]{Lang.java} ... \end{javacode}  (darf über Seiten umbrechen)
%
% Konsolenausgabe (Eingaben des Benutzers zwischen @...@):
%   \begin{konsole} Bitte Jahr eingeben:
%   @2026@ \end{konsole}
%
% Leere Konsole zum Ausfüllen (z. B. in Tests), Höhe in Zeilen:
%   \leerekonsole[5]
%%%%%%%%%%%%%%%%%%%%%%%%%%%%%%%%%%%%%%%%%%%%%%%%%%

\RequirePackage{xcolor}
\RequirePackage{listings}
\RequirePackage[most]{tcolorbox}
\RequirePackage[varqu,varl,scaled=0.95]{zi4}   % Inconsolata als Schreibmaschinenschrift

% Farben passend zu den Grafiken (Syntaxdiagramme, Dreieck, Pizza)
\definecolor{diagramm}{RGB}{74,144,217}
\definecolor{codeHintergrund}{RGB}{244,248,253}
\definecolor{codeSchluessel}{RGB}{28,86,160}
\definecolor{codeKommentar}{RGB}{125,132,145}
\definecolor{codeString}{RGB}{40,128,72}
\definecolor{codeZeilennr}{RGB}{150,170,195}
\definecolor{konsoleHintergrund}{RGB}{238,239,242}
\definecolor{konsoleText}{RGB}{45,50,60}
\definecolor{konsoleEingabe}{RGB}{204,102,0}
\definecolor{konsoleReiter}{RGB}{110,118,130}

% Umlaute in Kommentaren und Strings
\lstdefinestyle{gdiUmlaute}{
	literate=
		{ä}{{\"a}}1 {ö}{{\"o}}1 {ü}{{\"u}}1
		{Ä}{{\"A}}1 {Ö}{{\"O}}1 {Ü}{{\"U}}1
		{ß}{{\ss}}1
}

\lstdefinestyle{gdiJava}{
	style=gdiUmlaute,
	aboveskip=0pt,
	belowskip=0pt,
	language=Java,
	basicstyle=\ttfamily\footnotesize,
	keywordstyle=\color{codeSchluessel}\bfseries,
	commentstyle=\color{codeKommentar}\itshape,
	stringstyle=\color{codeString},
	numbers=left,
	numberstyle=\scriptsize\ttfamily\color{codeZeilennr},
	numbersep=10pt,
	tabsize=4,
	columns=fullflexible,
	keepspaces=true,
	showstringspaces=false,
	breaklines=true,
	postbreak=\space
}

\lstdefinestyle{gdiKonsole}{
	style=gdiUmlaute,
	aboveskip=0pt,
	belowskip=0pt,
	language={},
	basicstyle=\ttfamily\footnotesize\color{konsoleText},
	columns=fullflexible,
	keepspaces=true,
	breaklines=true,
	moredelim=[is][\color{konsoleEingabe}\bfseries]{@}{@}
}

% Gemeinsames Aussehen der Kästen
\tcbset{
	gdiKasten/.style={
		enhanced, breakable,
		listing only,
		sharp corners,
		boxrule=0pt,
		top=1.5mm, bottom=1.5mm, right=2mm,
		before skip=3mm, after skip=3mm,
		attach boxed title to top left={xshift=0mm, yshift=-0.2mm},
		boxed title style={sharp corners, boxrule=0pt, top=0.6mm, bottom=0.6mm, left=2mm, right=2mm},
		fonttitle=\ttfamily\footnotesize\bfseries
	}
}

% Quellcode: heller Kasten mit blauer Linie links und Dateiname als Reiter
\newtcblisting{javacode}[2][]{
	gdiKasten,
	listing options={style=gdiJava},
	colback=codeHintergrund,
	unbreakable,
	borderline west={2pt}{0pt}{diagramm},
	left=8mm,
	title={#2},
	coltitle=white,
	boxed title style={colback=diagramm},
	#1
}
% Kurzer Ausschnitt ohne Zeilennummern
\tcbset{Ausschnitt/.style={listing options={style=gdiJava, numbers=none}, left=3mm}}

% Konsolenausgabe: hellgrauer Kasten (druckerfreundlich), Eingaben orange
\newtcblisting{konsole}[1][Konsole]{
	gdiKasten,
	listing options={style=gdiKonsole},
	colback=konsoleHintergrund,
	unbreakable,
	left=3mm,
	title={#1},
	borderline west={2pt}{0pt}{konsoleReiter},
	coltitle=white,
	boxed title style={colback=konsoleReiter}
}

% Leere Konsole zum Ausfüllen (z. B. in Tests): \leerekonsole[Anzahl Zeilen]
\newcommand{\leerekonsole}[1][5]{%
	\begin{tcolorbox}[
		enhanced, unbreakable, sharp corners, boxrule=0pt,
		colback=konsoleHintergrund,
		borderline west={2pt}{0pt}{konsoleReiter},
		before skip=3mm, after skip=3mm,
		height=\dimexpr #1\dimexpr 8mm\relax + 3mm\relax,
		title={Konsole}, coltitle=white,
		attach boxed title to top left={xshift=0mm, yshift=-0.2mm},
		boxed title style={sharp corners, boxrule=0pt, colback=konsoleReiter, top=0.6mm, bottom=0.6mm, left=2mm, right=2mm},
		fonttitle=\ttfamily\footnotesize\bfseries
	]\end{tcolorbox}%
}
